\documentclass[aspectratio=169]{beamer}

% -------------------------------------------------------
% Lingua e pacchetti
\usepackage[italian]{babel}
\usepackage[T1]{fontenc}
\usepackage[utf8]{inputenc}
\usepackage{lmodern}
\usepackage{amsmath,amssymb,mathtools}
\usepackage{bm}

\title{Algoritmo GMRES}
\subtitle{Generalized Minimal Residual Method}
\author{}
\date{}

\begin{document}

% -------------------------------------------------------
\begin{frame}
  \titlepage
\end{frame}

% -------------------------------------------------------
\begin{frame}{Obiettivo del GMRES}
\textbf{Problema:} risolvere un sistema lineare di grandi dimensioni
\[
A x = b,
\]
dove $A \in \mathbb{R}^{n\times n}$ è in generale non simmetrica.

\medskip
\textbf{Idea:}
\begin{itemize}
  \item Costruire approssimazioni $x_k$ nello spazio di Krylov
  \[
  \mathcal{K}_k(A,r_0)=\mathrm{span}\{r_0, A r_0, A^2 r_0, \dots, A^{k-1} r_0\},
  \]
  dove $r_0=b-Ax_0$ è il residuo iniziale.
  \item Cercare $x_k$ che \emph{minimizza la norma del residuo}:
  \[
  \|r_k\| = \|b - A x_k\| = \min_{x \in x_0+\mathcal{K}_k(A,r_0)} \|b - A x\|.
  \]
\end{itemize}
\end{frame}

% -------------------------------------------------------
\begin{frame}{Forma del vettore di approssimazione}
Il vettore approssimato $x_k$ è della forma
\[
x_k = x_0 + V_k y_k,
\]
dove:
\begin{itemize}
  \item $V_k = [v_1, \dots, v_k]$ è la base ortonormale della sottospazio di Krylov (ottenuta con Arnoldi);
  \item $y_k \in \mathbb{R}^k$ è il vettore dei coefficienti da determinare;
  \item il residuo corrispondente è $r_k = b - A x_k$.
\end{itemize}

\medskip
Applicando la relazione di Arnoldi
\[
A V_k = V_{k+1}\,\underline{H}_k,
\]
si può scrivere
\[
r_k = \beta v_1 - V_{k+1}\underline{H}_k y_k,
\qquad \beta = \|r_0\|.
\]
\end{frame}

% -------------------------------------------------------
\begin{frame}{Problema ai minimi quadrati}
Poiché le colonne di $V_{k+1}$ sono ortonormali,
\[
\|r_k\| = \|\beta e_1 - \underline{H}_k y_k\|,
\]
dove $\underline{H}_k$ è una matrice $(k+1)\times k$ di Hessenberg superiore.

\medskip
\textbf{Problema ai minimi quadrati:}
\[
y_k = \arg\min_y \|\beta e_1 - \underline{H}_k y\|.
\]
\textbf{Soluzione:}
\begin{itemize}
  \item Il problema è di dimensione piccola ($k\times k$);
  \item Si risolve facilmente con fattorizzazione QR o rotazioni di Givens;
  \item Il residuo minimo è $\|r_k\| = \|\beta e_1 - \underline{H}_k y_k\|$.
\end{itemize}
\end{frame}

% -------------------------------------------------------
\begin{frame}{Algoritmo GMRES passo per passo}
\textbf{Input:} $A$, $b$, $x_0$, tolleranza, numero massimo di passi $m$.

\medskip
\textbf{1.} Calcola $r_0 = b - A x_0$, $\beta = \|r_0\|$, $v_1 = r_0 / \beta$.

\textbf{2.} Per $k = 1,2,\dots,m$:
\begin{enumerate}
  \item Esegui \textbf{una} iterazione di Arnoldi per ottenere $v_{k+1}$ e $\underline{H}_k$;
  \item Aggiorna il problema ai minimi quadrati
        $\min_y \|\beta e_1 - \underline{H}_k y\|$
        usando rotazioni di Givens;
  \item Se $\|r_k\|/\|b\|$ è sotto la tolleranza, \textbf{stop}.
\end{enumerate}

\textbf{3.} Calcola $y_k$ (back-substitution) e
\[
x_k = x_0 + V_k y_k.
\]
\end{frame}

% -------------------------------------------------------
\begin{frame}{Caratteristiche principali}
\begin{itemize}
  \item Il metodo \textbf{non costruisce} $x_{k+1}$ da $x_k$ direttamente,
        ma trova la migliore approssimazione nello spazio $\mathcal{K}_k(A,r_0)$.
  \item Il processo è \textbf{incrementale}:
        ad ogni passo si aggiunge un nuovo vettore alla base di Arnoldi.
  \item Il residuo $\|r_k\|$ è \textbf{monotonicamente decrescente}.
  \item È necessaria solo la moltiplicazione $A v_k$ (non l’inversione di $A$).
  \item Se la memoria cresce troppo, si usa la versione \textbf{GMRES(m)} con \emph{restarts}.
\end{itemize}
\end{frame}

% -------------------------------------------------------
\begin{frame}{GMRES con restart}
\textbf{Problema:} il costo e la memoria crescono con il numero di passi $k$ (serve salvare tutti i $v_i$).

\medskip
\textbf{Soluzione: Restart GMRES (GMRES(m)):}
\begin{enumerate}
  \item Esegui GMRES per $m$ passi $\Rightarrow$ ottieni $x_m$;
  \item Usa $x_0 \leftarrow x_m$ come nuova approssimazione iniziale;
  \item Ricomincia il processo con il nuovo residuo $r_0=b-Ax_0$.
\end{enumerate}

\textbf{Osservazione:} GMRES(m) mantiene la memoria costante ma può convergere più lentamente.
\end{frame}

% -------------------------------------------------------
\begin{frame}{Riepilogo visivo}
\begin{itemize}
  \item GMRES costruisce una successione di sottospazi di Krylov:
  \[
  \mathcal{K}_1 \subset \mathcal{K}_2 \subset \dots \subset \mathcal{K}_k.
  \]
  \item Ad ogni passo:
  \begin{enumerate}
    \item Si espande la base ortonormale (Arnoldi);
    \item Si risolve un piccolo problema ai minimi quadrati;
    \item Si ottiene la soluzione $x_k$ che minimizza $\|r_k\|$.
  \end{enumerate}
  \item La norma del residuo decresce e la soluzione converge a quella esatta.
\end{itemize}
\end{frame}

% -------------------------------------------------------
\begin{frame}{Vantaggi e limiti del GMRES}
\textbf{Vantaggi:}
\begin{itemize}
  \item Funziona per matrici non simmetriche;
  \item Minimizza esplicitamente la norma del residuo;
  \item Non richiede la conoscenza della matrice inversa.
\end{itemize}

\textbf{Limiti:}
\begin{itemize}
  \item Crescita del costo e della memoria con $k$;
  \item Richiede ortogonalizzazione di tutti i vettori precedenti;
  \item Può perdere ortogonalità numerica (richiede ri-ortogonalizzazione);
  \item Possibile convergenza lenta senza \textbf{precondizionamento}.
\end{itemize}
\end{frame}

% -------------------------------------------------------
\begin{frame}{Schema sintetico dell'algoritmo GMRES}
\begin{block}{}
\begin{align*}
r_0 &= b - A x_0, \qquad \beta = \|r_0\|, \quad v_1 = r_0 / \beta,\\
\text{per } j &= 1,2,\dots,k:\\
&\quad w = A v_j,\\
&\quad \text{per } i=1,\dots,j:\; h_{ij} = v_i^\top w,\; w \leftarrow w - h_{ij}v_i,\\
&\quad h_{j+1,j} = \|w\|,\; v_{j+1} = w / h_{j+1,j},\\
&\quad \text{aggiorna } \|\beta e_1 - \underline{H}_j y\| \text{ con rotazioni di Givens.}
\end{align*}
\end{block}

\[
x_k = x_0 + V_k y_k, \qquad \|r_k\| = \|\beta e_1 - \underline{H}_k y_k\|.
\]
\end{frame}

\end{document}

